\section*{Archived Reader's Guide and working roadmap from the public v3.13.0 release}
\addcontentsline{toc}{section}{Archived Reader's Guide and working roadmap from the public v3.13.0 release}

The following note is preserved from the public v3.13.0 release stage. It records the release-facing reader's guide and roadmap exactly for that checkpoint; in the present release continuation it is archival proof-memory rather than a current status block.
It presents the existing mathematical line in one monolithic text, preserves the stabilized diagnostic material already carried by the document, and keeps the operative rule:
\emph{do not shorten, do not overwrite, only sharpen, normalize, and preserve.}

\subsection*{Archived version identity for the public v3.13.0 checkpoint}
\begin{itemize}[leftmargin=2em]
  \item Historical version: 3.13.0
  \item Historical date: 2026-04-07
  \item Concept DOI: \href{https://doi.org/10.5281/zenodo.19210728}{10.5281/zenodo.19210728}
  \item Historical status: public Zenodo release v3.13.0
  \item Historical public release DOI (v3.13.0): \href{https://doi.org/10.5281/zenodo.19446872}{10.5281/zenodo.19446872}
  \item Historical version-specific DOI for v3.13.0: \href{https://doi.org/10.5281/zenodo.19446872}{10.5281/zenodo.19446872}
  \item Previous Zenodo releases: v3.3.5 (\href{https://doi.org/10.5281/zenodo.19396143}{10.5281/zenodo.19396143}), v3.4.0 (\href{https://doi.org/10.5281/zenodo.19410949}{10.5281/zenodo.19410949}), v3.7.1 (\href{https://doi.org/10.5281/zenodo.19412478}{10.5281/zenodo.19412478}), v3.9.0 (\href{https://doi.org/10.5281/zenodo.19423689}{10.5281/zenodo.19423689})
  \item Main-paper relation: Explicit mathematical companion to the current Sphaera master line~\cite{brendecke2026main}
\end{itemize}

\subsection*{Archived working note from v3.13.0}
This public release proceeds from the public v3.11.0 Companion release in one preserved monolithic text. Content-wise it now keeps together three central gains:
\begin{enumerate}[leftmargin=2em,label=(\roman*)]
  \item the stabilized residual-record interface and its associated structural tools,
  \item the reproducible M\"obius-shell diagnostic suite for the live residual scales, and
  \item the OP~5 packet, breathing-envelope, and frame-flicker visual audit line.
\end{enumerate}
Accordingly, that archived public release should be read as the stabilized public checkpoint of its then-current monolithic chain: it preserved the mathematical line of that stage while keeping the front cleaner and the live burdens more sharply routed.

\subsection*{Archived primary working goals recorded at v3.13.0}
\begin{enumerate}[leftmargin=2em]
  \item Present the current Companion mathematics in one public monolithic release.
  \item Keep the stabilized generator--atlas--holonomy core and the residual-record package $\mathfrak R(B)=\bigl(\Phi,\Delta_{\mathrm{diag}},\varepsilon_{\mathrm{par}},\chi_{\mathrm{par}},\widehat\chi_{\mathrm{par}}\bigr)$ intact as the faithful visible interface.
  \item Preserve the M\"obius-shell diagnostics as reproducible, active-PDF figures so that residual drift, continuation, sensitivity, and sector persistence remain directly auditable.
  \item Preserve the OP~5 packet-wave, breathing-envelope, and frame-flicker visual line in one release-safe appendix block.
  \item Preserve the monolithic document without shortening historical material.
\end{enumerate}

\subsection*{Archived release priorities recorded at v3.13.0 toward Zenodo deposition}
\begin{enumerate}[leftmargin=2em]
  \item Turn the current OP~5 packet/phase/epsilon analysis into one cleaner structural branch theorem with a small admissible branch load and a non-escape criterion.
  \item Keep OP~1 on the transported-branch reading and compress the remaining burden to the discrete correction channel without reopening settled continuum structure.
  \item Normalize the document language so that internal patch chronology remains archival only, while the main line reads as one coherent proof route.
  \item Collect the now-stabilized transport, readout, checkpoint, and epsilon-trap principles as explicit reusable tools for the next theorem blocks.
  \item Use the present v3.13.0 public release as the stabilized next checkpoint because the OP transport closure architecture is now integrated more consistently across the monolithic text rather than remaining split between new theorem blocks and unreconciled historical passages.
\end{enumerate}

\subsection*{Continuity of the v3.5.2 checkpoint in the present release}
The present file preserves the former \emph{v3.5.2 abstract-polish checkpoint} as compatible earlier material within the current release.
Concretely, the 3.5.2 achievements remain part of the present monolith:
\begin{enumerate}[leftmargin=2em,label=(\roman*)]
  \item the residual-record interface \(\mathfrak R(B)\) and the reconstruction normal form remain unchanged;
  \item the readout-relative lemma methodology, structural lift logic, and Theorem~7 / Hypotheses~7.4--7.6 remain preserved as dependencies;
  \item the former OP~2 operator burden is no longer carried as a separate primary front, but is preserved in reclassified form inside the later OP-first status reading;
  \item the abstract-polish framing of 3.5.2 is therefore not discarded, but sharpened and transported into the later residual-controlled and boundary-layer language.
\end{enumerate}
Accordingly, the development line may be read as:
\[
\begin{aligned}
&\text{v3.5.2 interface normalization}
\Longrightarrow
\text{v3.6 conservative frontend}
\Longrightarrow
\text{v3.7 OP-first stabilization}\\[0.2em]
&\Longrightarrow
\text{v3.8 residual / boundary-layer sharpening}.
\end{aligned}
\]


\subsection*{Archived current mathematical status at the public v3.13.0 release stage}
For the current release line, the reader should keep one status separation in mind before entering the preserved historical layers below:
\begin{itemize}[leftmargin=2em]
  \item The stabilized generator--atlas--holonomy sector from v3.4 is already closed at the structural level and has been compressed into the admissible residual record
  \[
  \mathfrak R(B)=\bigl(\Phi,\Delta_{\mathrm{diag}},\varepsilon_{\mathrm{par}},\chi_{\mathrm{par}},\widehat\chi_{\mathrm{par}}\bigr).
  \]
  \item The current v3.5.x step is therefore no longer a search for new visible primitives, but a normalization of comparison, transport, and interpretation through that record.
  \item Historical OP~2 should no longer be read as a separate live operator problem; its T1/T2/T3 blocks are closed, and the former T4 frame-offset part is absorbed into the residual quantitative completion of OP~1.
  \item Historical OP~3 is closed at the present structural level; it may still be enriched later, but it is no longer one of the primary live fronts.
  \item Historical OP~5 has been reduced to the photon-frame Step~(B) remainder only; Step~(A) is already closed in the present text.
  \item The genuinely live fronts are therefore the residual continuum-to-measurement correction for \(\bsym\), controlled emergence with explicit constants and bounds, the quantitative mass hierarchy, the three-generation splitting law, the compact Dixon-algebra reformulation, and the remaining Theta-sum cancellation step.
\end{itemize}

\paragraph{Release-stage caution.}
The present reading is strong enough for this release, but it is intentionally stated in stabilized form: the remaining live fronts are read as residual downstream tasks \emph{in the present stabilized reading}. This is stronger than a loose hypothesis, because it rests on the exact representability and conservative-frontend results of Sections~3.5--3.6; but it is still a mathematically revisable reading if a later residual OP were to exhibit a genuinely missing visible degree of freedom. No such counter-indication is currently visible in the Companion chain.

\paragraph{How to read the preserved historical roadmap blocks.}
The longer roadmap sections that follow from earlier cycles are retained deliberately as proof-memory and dependency control. They should now be read as \emph{historical support strata}: useful for reconstruction, closure auditing, and tool extraction, but no longer identical with the short list of primary live questions at the head of the present line.

\subsection*{What the 3.5.x line achieved, and what 3.6 now fixes first}
At this point the 3.5.x line should be read as a \emph{completed interface-normalization stage} rather than as an unfinished search for further visible structure.
Its concrete achievements are:
\begin{enumerate}[leftmargin=2em]
  \item the full 3.4 generator--atlas--holonomy package has been compressed into the admissible residual record
  \[
  \mathfrak R(B)=\bigl(\Phi,\Delta_{\mathrm{diag}},\varepsilon_{\mathrm{par}},\chi_{\mathrm{par}},\widehat\chi_{\mathrm{par}}\bigr);
  \]
  \item visible reconstruction has been normalized through the record and the chartwise shell data \((s,\pi)\);
  \item record-level comparison and transport have been fixed through admissible morphisms and residual-record isomorphisms;
  \item the faithful interpreter principle has made the stabilized visible sector exactly representable on record level;
  \item the historical OP-layer has been reclassified so that closed, reduced, and genuinely open fronts are no longer mixed.
\end{enumerate}
Accordingly, the next step should \emph{not} attempt to reopen the visible interface.
A mathematically clean 3.6 target is instead:
\begin{enumerate}[leftmargin=2em,label=(\roman*)]
  \item prove the conservative-frontend principle explicitly, i.e. every later language layer must factor through the residual-record interpreter and may not introduce any new visible primitive;
  \item sharpen the remaining live fronts only as downstream quantitative tasks: prime-lattice/continuum correction for \(\bsym\), explicit constants and bounds for controlled emergence, mass hierarchy and three-generation splitting, compact Dixon-algebra reformulation, and the unresolved Theta-sum cancellation block;
  \item prepare the eventual ZFM-/language-layer only after the conservative factorization theorem is in place.
\end{enumerate}
In short: \emph{3.5.x stabilized the interface; 3.6 begins the stabilization of the admissible languages above it.}

\subsection*{Reading discipline for the next stage}
Every new subsection introduced below is to be read under four status labels that must be kept distinct:
\begin{enumerate}[leftmargin=2em,label=(\alph*)]
  \item structurally forced,
  \item projectively read,
  \item interpretive overlay separated from theorem status,
  \item reserved for later proof expansion.
\end{enumerate}


\subsection*{Archived mathematics-first completion path realized in v2.27.0}

The present major release is to be read as a \emph{mathematics-first completion cycle}. The computational chapter opened in v2.24.0--v2.26.0 is to remain in place as a visible teaser and reserved architecture, but it is no longer the primary work driver for the immediate next version. For the next cycle, the operative rule is instead:
\begin{center}
\emph{close the mathematics first; let the proof engine wait for a sufficiently closed mathematical layer.}
\end{center}

Accordingly, the completion path realized in v2.27.0 is organized into five mathematical work packages:
\begin{enumerate}[leftmargin=2em,label=\textbf{WP\arabic*:}]
  \item \textbf{Core structural closure.} Complete and normalize the remaining structural chains around ground condition, HC emergence, closure admissibility, and the Millennium matrix as global HC realization.
  \item \textbf{Proof-chain hardening.} Close missing lemma-to-theorem dependencies, remove duplicated or merge-residual argument fragments, and restore explicit proof-memory where shortened passages still rely on hidden expansion.
  \item \textbf{Projector / transport mathematics.} Systematize projector logic, carrier-group-readout transport, and the semantic/tensorial kernel interface so that projection steps are mathematically explicit before they are computationally compressed.
  \item \textbf{Tools as mathematical appendices.} Gather structural tools, computational shortcuts, closure tricks, and trace/branch transfer rules as named mathematical instruments with scope conditions, rather than letting them remain scattered as implicit know-how.
  \item \textbf{Computational strand on hold-but-present.} Retain MML, MMA, MCP, certificates, and host export as a prepared future chapter, but do not let that layer outrun the mathematical substance it is supposed to compress.
\end{enumerate}

In practical terms, v2.27.0 therefore prioritizes mathematical completion over computational expansion. The computational chapter remains part of the Companion and should be preserved, cited, and gently normalized where needed, but only as a secondary architecture until the mathematical layer is sufficiently closed. The version-specific DOI for the present v2.27.0 release is \href{https://doi.org/10.5281/zenodo.19266666}{10.5281/zenodo.19266666}.


\subsection*{Operational mathematical closure agenda}

To make the mathematics-first rule operational rather than merely declarative, the next revision cycle is to be executed along the following closure agenda:
\begin{enumerate}[leftmargin=2em,label=\textbf{MA\arabic*:}]
  \item \textbf{Definition pass.} Re-read and normalize the basic definitions that the later chains depend on: collapse state, closure, admissibility, projector sector, compensator action, transport layer, carrier/group/readout distinction, and Millennium-matrix status.
  \item \textbf{Dependency pass.} Identify theorem blocks whose proofs still rely on compressed or implicit dependencies, then make the lemma chain explicit enough that each forward step can be reconstructed without external memory.
  \item \textbf{Transport pass.} Isolate the mathematics of projection, frame transport, Lorentz-mediated readout, and disjoint-sector return so that these are expressed first as mathematical maps and only later as computational metaphors.
  \item \textbf{Appendix-tools pass.} Move recurrent structural tricks into named mathematical tools with hypotheses and scope, so that later proof reuse becomes mathematically transparent rather than conversationally remembered.
  \item \textbf{Stability pass.} Check the normalized chapter sequence against merge-residual duplication, outdated internal cross-references, and proof steps that were historically preserved but not yet fully re-hardened.
\end{enumerate}

\subsection*{Priority clusters for v2.27.0}

The following clusters should receive mathematical priority before any further major proof-engine expansion:
\begin{enumerate}[leftmargin=2em,label=\textbf{PC\arabic*:}]
  \item \textbf{Ground condition $\Rightarrow$ HC $\Rightarrow$ confinement.} This chain should be rewritten as a clean mathematical dependency, with the Heisenberg compensator treated as emergent from the ground condition rather than as an independently floating postulate.
  \item \textbf{Projector and closure mathematics.} The admissible-sector logic should be made fully explicit, including which projectors are primitive, which are derived, and how closure is checked at the mathematical rather than merely intuitive level.
  \item \textbf{Millennium matrix as mathematical object.} The Millennium matrix should first be stabilized as a global mathematical realization of the HC/closure architecture before any stronger computational reading is allowed to dominate the exposition.
  \item \textbf{Trace / branch mathematics.} Probe traces, branch transfer, Hopf-transition logic, and return structure should be stated as mathematical transformations with hypotheses and admissibility conditions, not only as algorithmic motifs.
  \item \textbf{Proof-memory restoration.} Where earlier versions compressed long chains into structural shorthand, the next cycle should restore enough explicit mathematics that the shorthand remains auditable and publication-safe.
\end{enumerate}

\subsection*{What is deliberately deferred beyond v2.27.0}

The following lines are explicitly preserved but not treated as primary completion goals for the next release cycle:
\begin{enumerate}[leftmargin=2em,label=\textbf{D\arabic*:}]
  \item full MML language formalization,
  \item full MMA/MCP execution refinement,
  \item standalone proof-kernel extraction,
  \item stronger universality claims,
  \item operating-system style runtime architecture.
\end{enumerate}
These elements remain valuable, but in the current sequencing they are to follow mathematical closure rather than compete with it.

\subsection*{Mathematical work matrix for the next pass}
To make the mathematics-first agenda operational, the next pass is organized as a compact work matrix rather than as a loose wish list.
\begin{center}
{\scriptsize\setlength{\tabcolsep}{2pt}\begin{tabular}{@{}p{0.18\textwidth}p{0.24\textwidth}p{0.27\textwidth}p{0.06\textwidth}@{}}
\toprule
\textbf{Block} & \textbf{Target state} & \textbf{Main dependencies} & \textbf{Priority}\\
\midrule
Ground condition $\to$ HC $\to$ confinement & explicit statement chain, normalized notation, and no ambiguity about emergence vs. axiomatics & core structural assumptions and septinion confinement language & A\\
Projector and closure mathematics & sharpened lemmas, admissibility conditions, and stable projector notation & closure operator layer, readout discipline, sector decomposition & A\\
Millennium matrix as mathematical object & operator-level formulation, domain/codomain clarity, and clarified relation to HC & projector layer, transport layer, compensator interpretation & A\\
Trace / branch / Hopf transport mathematics & clean separation of trace, branching, return, and transport rules & probe framework, disjointness, Lorentz-mediated transport & B\\
Dependency chain hardening & theorem/lemma dependency map cleaned and checked for merge artifacts & current theorem chain, proposition and lemma references & A\\
Structural tools appendix & collect reusable mathematical shortcuts without shortening proofs & established proof moves, normalized notation & B\\
Proof-memory restoration / archival recovery & recover or reintegrate historically displaced proof material where needed & older Companion states, recovered large basis, appendix chronology & B\\
Computational teaser retention & keep chapter visible, but frozen except for compatibility edits & mathematical closure progress & C\\
\bottomrule
\end{tabular}}
\end{center}

A-priority blocks should drive the next version boundary. B-priority blocks should be developed in service of those A-blocks. C-priority blocks should be touched only when required for compatibility with the mathematical reorganization.

\subsection*{Dependency map of the mathematical core}
To prevent the next pass from drifting into local edits without structural order, the mathematical core is also organized as a dependency map. The intended reading is simple: if an upstream block remains unstable, all downstream blocks that consume it can only be provisionally edited.
\begin{center}
\begin{tabular}{p{0.28\textwidth}p{0.30\textwidth}p{0.30\textwidth}}
\toprule
\textbf{Core block} & \textbf{Direct prerequisite} & \textbf{Main downstream consumers}\\
\midrule
Ground condition and collapse picture & none; foundational starting point & HC emergence, confinement thesis, collapse-state language, admissibility reading\\
HC emergence and confinement chain & ground condition, collapse picture & projector logic, trace/branch admissibility, Millennium-matrix interpretation\\
Projector / closure mathematics & HC chain, admissibility normalization & branch filtering, Boolean readout, proof templates, certificate layer\\
Transport / Lorentz layer & collapse picture, projector sector discipline & return-map reading, branch transport, reflector-like structural reversals\\
Millennium matrix as mathematical object & HC chain, projector/closure layer, transport layer & computational teaser chapter, co-processor view, host-export discipline\\
Trace / branch mathematics & projector logic, transport layer, admissibility conventions & Hopf branching examples, trace reconciliation, branch templates\\
Tool and appendix lemmas & all stabilized upstream definitions & proof hardening, auditability, future proof-engine extraction\\
\bottomrule
\end{tabular}
\end{center}

In operational terms, this means that the next mathematical pass should proceed roughly in the order
\[
\begin{aligned}
&(\text{ground condition} \Rightarrow \text{HC} \Rightarrow \text{confinement}) \prec (\text{projector/closure})\\
&\prec (\text{transport}) \prec (\text{Millennium matrix, trace, branch})\\
&\prec (\text{appendix tools and supporting lemmas}).
\end{aligned}
\]
Any local repair that violates this order may still be useful as a note, but it should not be mistaken for structural closure.


\noindent\textbf{Definition-and-notation pass for the next mathematical stage.} Before the next major release, the Companion should also undergo a focused definition-and-notation pass. The purpose of this pass is not to shorten the document, but to reduce merge drift, ambiguous symbol reuse, and status confusion between theorem, lemma, tool, hypothesis, and teaser language.

\paragraph{Mathematically fixed for the current line.}
For the current Companion line, the following should be treated as fixed structural anchors unless explicitly revised in a later release: the collapse-state viewpoint as the local readout picture; the phase-shifted measurement-frame perspective; the ground-condition $\Rightarrow$ Heisenberg-compensator $\Rightarrow$ confinement chain; the projector/closure discipline for admissible sectors; the standing triolen architecture with one orthogonal component; and the imaginary convention already adopted for the relevant octonionic/triadic units.

\paragraph{Notation to normalize before v2.27.0.}
Notation should be normalized especially in the following zones: projector symbols and their admissibility role; closure notation versus stabilization notation; transport/Lorentz notation across layers; Millennium-matrix notation when used as a mathematical object rather than as a computational teaser; trace versus probe-trace notation; branch/Hopf notation; and the distinction between theorem-level mathematical statements and later computational analogies. In particular, the same symbol should not silently alternate between mathematical object, physical readout, and computational metaphor.

\paragraph{Items that must remain hypotheses.}
Some parts of the Companion should remain explicitly flagged as hypotheses, interpretive proposals, or deferred claims rather than being silently upgraded into theorems. This applies in particular to the Higgs/coupling interpretation, to stronger black-hole or event-horizon extrapolations, to stronger computational-universality language, and to any claim that would identify the computational chapter with a finished proof engine rather than a prepared architecture. The next mathematical pass should therefore favor closure and definition-hardening over conceptual inflation.

\paragraph{Operational rule for the next pass.}
For the next revision cycle, each touched block should be tagged mentally by the question: \emph{definition, theorem, lemma, tool, hypothesis, or teaser?} The next major mathematical gain is expected to come more from making these statuses explicit and stable across the Companion than from adding further new conceptual branches.


\subsection*{Target-state matrix for v2.27.0}

\small\begin{longtable}{p{0.22\textwidth}p{0.17\textwidth}p{0.24\textwidth}p{0.25\textwidth}}
\toprule
\textbf{Core block} & \textbf{Current status} & \textbf{Target state for v2.27.0} & \textbf{Minimal completion condition} \\
\midrule
\endfirsthead
\toprule
\textbf{Core block} & \textbf{Current status} & \textbf{Target state for v2.27.0} & \textbf{Minimal completion condition} \\
\midrule
\endhead
Ground condition / collapse picture & Structurally central, but still partly distributed across explanatory passages. & Gathered into a visibly controlling mathematical block. & The collapse-state picture is stated once in stable mathematical language and used consistently as the local readout frame. \\

HC emergence / confinement chain & Present as a guiding structural principle, but not yet fully normalized as a dependency spine. & Explicitly treated as a mathematical chain: ground condition $\Rightarrow$ HC $\Rightarrow$ confinement. & The chain appears in one normalized theorem/lemma sequence and is referenced consistently downstream. \\

Projector / closure mathematics & Strongly present, but notation and admissibility wording still need tightening. & Treated as the first closed admissibility layer of the framework. & Projector symbols, closure symbols, and admissibility language are normalized and no longer drift between sections. \\

Transport / Lorentz layer & Conceptually active, but still mixed between mathematical transport and physical interpretation. & Clarified as the transport layer relating frames without overloading its physical extrapolations. & A stable mathematical transport notation is fixed and separated from speculative physical remarks. \\

Millennium matrix as mathematical object & Partly present beside its computational reading. & Brought back under mathematical control before any stronger proof-engine use. & The matrix is stated first as a mathematical object/operator, with the computational reading clearly subordinate. \\

Trace / branch mathematics & Available in structural language, but still partly entangled with later computational analogies. & Recast as a mathematical trace/branch layer before proof-engine expansion resumes. & Trace, probe-trace, branch, and Hopf language are normalized and tied to mathematical dependence rather than metaphor. \\

Tool / appendix lemmas & Growing, but not yet fully separated from main-line dependencies. & Reorganized as supporting tools rather than semi-primary arguments. & Each tool is clearly marked as tool/appendix support and no longer silently carries theorem-level load without declaration. \\
\bottomrule
\end{longtable}

\paragraph{Use of the matrix.}
This target-state matrix is intended as an operational filter for the next revision cycle. A block should only be counted as \emph{mathematically advanced} if its current status has moved measurably toward the declared target state and if its minimal completion condition has been checked explicitly.

\subsection*{Proof-dependency checklist for v2.27.0}
To keep the next mathematical pass referentially stable, the following dependency checklist should be treated as a minimal proof-control layer. The point is not to produce exhaustive dependency graphs for every sentence, but to ensure that the main mathematical blocks are supported by a sufficiently explicit chain of definitions, lemmas, transport statements, and appendix tools.

\begin{itemize}
    \item \textbf{Ground condition / collapse picture.} Must rest on stable base definitions, a clear distinction between structural state and measurement readout, and an explicit statement of which consequences are structural and which remain interpretive.
    \item \textbf{HC emergence / confinement chain.} Must depend only on the ground condition and the confinement logic, not on later computational analogies. For v2.27.0 the chain ground condition $\Rightarrow$ HC $\Rightarrow$ confinement should be readable without hidden appeals to deferred sections.
    \item \textbf{Projector / closure mathematics.} Requires normalized projector notation, at least one closure lemma family, and stable cross-references into the appendix tools where routine projector identities are collected.
    \item \textbf{Transport / Lorentz layer.} Must state exactly which structural transformations are taken as transport rules, which are theorem-level, and which remain hypothesis-level interpretations. Dependencies on projector compatibility should be explicit.
    \item \textbf{Millennium matrix as mathematical object.} Should depend on the already stabilized projector/closure and transport layer, rather than carrying foundational burden by itself. For v2.27.0 the matrix should be mathematically positioned after the base closure apparatus, not before it.
    \item \textbf{Trace / branch mathematics.} Requires a small stable family of supporting lemmas on admissible continuation, branch exhaustion, and trace comparison. The next pass should avoid leaving branch claims supported only by prose.
    \item \textbf{Appendix tools and supporting lemmas.} These should be treated as genuine support infrastructure: enough to stabilize the main line, but not so invasive that the main thread becomes unreadable. The key requirement is reference stability and not merely accumulation.
\end{itemize}

For the purpose of the next release threshold, a proof block should count as \emph{referentially stable} only if its direct supporting definitions and lemmas can be named without ambiguity, found in the current Companion without version drift, and read in an order compatible with the mathematical dependency map already given above.


\subsection*{Reference-stability pass for v2.27.0}

For the mathematics-first pass toward v2.27.0, internal references are themselves treated as part of structural closure. The next public version should not only add arguments, but also stabilize the minimal web of labels, theorem references, and section anchors on which the A-priority mathematical blocks depend.

\begin{itemize}
  \item \textbf{Definition stability.} The core labels for the ground condition, Heisenberg-compensator emergence, confinement statement, projector/closure layer, transport layer, and Millennium-matrix mathematical section should remain fixed under the next hardening pass unless an explicit renumbering note is added.
  \item \textbf{Dependency stability.} Every A-priority theorem or proposition should point only to labels that are themselves already fixed at the level of definition, lemma, theorem, or supporting tool. Soft roadmap anchors should not be used as silent proof dependencies.
  \item \textbf{Appendix-tool stability.} If a proof uses an appendix tool, the tool must be referenced by a stable label and its status must remain visibly marked as tool, lemma, or hypothesis.
  \item \textbf{Hypothesis transparency.} Any block that still depends on a hypothesis label must keep that hypothesis status explicit in the surrounding prose and may not inherit theorem-strength merely by cross-reference.
  \item \textbf{Public-release threshold.} For v2.27.0, a block counts as reference-stable only if its internal labels, backward references, and appendix-tool dependencies survive a clean compilation pass without ad hoc renaming.
\end{itemize}

Operationally, the next mathematical stage should therefore include a short explicit reference-stability check before public release: (i) labels on A-priority definitions and theorems are frozen, (ii) proof dependencies resolve only to mathematically declared objects, (iii) appendix tools referenced in the main line remain stable and typed, and (iv) theorem-to-hypothesis drift is not introduced by accidental cross-linking.

\subsection*{Block-by-block closure checklist for the next pass}
For each A-priority block, the next pass should explicitly separate three statuses: what remains to be closed, what may count as done for v2.27.0, and what may remain open without blocking the release threshold.

\begin{itemize}
    \item \textbf{Ground condition / collapse picture.} Missing: final normalized wording and explicit distinction between structural and observational language. Done for v2.27.0: stable definition and one clear theorem-level consequence. May remain open: broader interpretive extensions.
    \item \textbf{HC emergence / confinement chain.} Missing: explicit derivational chain and at least one supporting lemma that fixes the confinement reading. Done for v2.27.0: readable chain ground condition $\Rightarrow$ HC $\Rightarrow$ confinement. May remain open: stronger physical extrapolations.
    \item \textbf{Projector / closure mathematics.} Missing: normalized notation and a compact closure lemma family. Done for v2.27.0: a stable projector/closure core with appendix support. May remain open: larger auxiliary projector calculus.
    \item \textbf{Transport / Lorentz layer.} Missing: clean split between structural transport statements and hypothesis-level interpretive claims. Done for v2.27.0: one stable transport layer attached to the projector/closure block. May remain open: more speculative extensions.
    \item \textbf{Millennium matrix as mathematical object.} Missing: definitive positioning relative to HC, projector structure, and transport. Done for v2.27.0: mathematically explicit role after closure/transport stabilization. May remain open: the computational reading beyond teaser level.
    \item \textbf{Trace / branch mathematics.} Missing: compact lemma support for admissible continuation, exhaustion, and trace comparison. Done for v2.27.0: one stable branch/trace family sufficient for the main line. May remain open: extended proof-program analogies.
    \item \textbf{Appendix tools / supporting lemmas.} Missing: sharper separation between indispensable support and purely archival material. Done for v2.27.0: tool appendix that stabilizes references without overgrowing the main line. May remain open: secondary tool families.
\end{itemize}

\subsection*{Explicit subversion discipline}

From this point onward, intermediate work states are to be recorded explicitly as local subversions whenever substantial planning, normalization, recovery, or integration work has been performed without yet justifying a new Zenodo release. The intended pattern is:
\begin{center}
\texttt{public release v2.27.0 -> local subversion v2.27.1 -> further local subversions as needed -> next public milestone}
\end{center}
This rule is not cosmetic. It is meant to prevent silent rollback, silent overwrite, and accidental compression of the project memory inside a single file name.


\subsection*{Concrete chapter scaffold to be unfolded in the next revision cycle}

\paragraph{Part I --- Structural Core}
\begin{enumerate}[leftmargin=2em,label=\arabic*.]
  \item \textbf{Introduction}
    \begin{enumerate}[label*=\arabic*., leftmargin=2.5em]
      \item Motivation
      \item Scope of this version
      \item Relation to the main paper
      \item Structural strategy: structure first, closure before projection
    \end{enumerate}
  \item \textbf{Structural Preliminaries}
    \begin{enumerate}[label*=\arabic*., leftmargin=2.5em]
      \item Basic structural units
      \item Orthogonality constraint
      \item Complex and imaginary structural setting
      \item Frame shift of measurement
      \item Structural closure intuition
    \end{enumerate}
  \item \textbf{Ground Condition and Emergent Closure}
    \begin{enumerate}[label*=\arabic*., leftmargin=2.5em]
      \item Ground condition
      \item From ground condition to HC
      \item HC as compensating constraint
      \item From HC to closure
      \item Protected invariants
      \item Proposition block: Ground Condition $\Rightarrow$ HC $\Rightarrow$ Closure
    \end{enumerate}
  \item \textbf{Septinion / Seption Structural Embedding}
    \begin{enumerate}[label*=\arabic*., leftmargin=2.5em]
      \item From triadic units to higher structural coupling
      \item Octonionic and higher assembly logic
      \item Septinion / seption confinement
      \item Structural consequences
    \end{enumerate}
  \item \textbf{Structural Master Equation --- First Layer}
    \begin{enumerate}[label*=\arabic*., leftmargin=2.5em]
      \item Why a master equation emerges
      \item Pure structural form
      \item Term-by-term structural origin
      \item Structural vs.\ projective terms
      \item First consistency remarks
    \end{enumerate}
\end{enumerate}

\paragraph{Part II --- Probe Sector}
\begin{enumerate}[leftmargin=2em,label=\arabic*.]
  \setcounter{enumi}{5}
  \item \textbf{Probe Geometry}
    \begin{enumerate}[label*=\arabic*., leftmargin=2.5em]
      \item Two probes moving toward each other
      \item Each probe as a triolic structure
      \item Orthogonal third group
      \item Definition of the search space
    \end{enumerate}
  \item \textbf{Probe Dynamics and Field Structure}
    \begin{enumerate}[label*=\arabic*., leftmargin=2.5em]
      \item Structural interaction regime
      \item Closure in the probe sector
      \item Probe sector equation
      \item Search, trace, and structural read-off
    \end{enumerate}
  \item \textbf{Probe Tools and Guided Search}
    \begin{enumerate}[label*=\arabic*., leftmargin=2.5em]
      \item Frequency and pattern detection
      \item Routing and candidate compression
      \item Transition guidance
      \item Convergence tools
      \item Optional statistical layer
      \item Interpretation as structural toolkit
    \end{enumerate}
\end{enumerate}

\paragraph{Part III --- Technical Normalization}
\begin{enumerate}[leftmargin=2em,label=\arabic*.]
  \setcounter{enumi}{8}
  \item \textbf{Structural Tools and Computational Shortcuts}
    \begin{enumerate}[label*=\arabic*., leftmargin=2.5em]
      \item Why tool normalization is needed
      \item Closure-tools
      \item Frame-shift-tools
      \item Projection-tools
      \item Probe-tools
      \item Equivalence-tools
      \item Read-off / trace-tools
      \item Standard form for every tool
    \end{enumerate}
\end{enumerate}

\paragraph{Part IV --- Interpretive Bridge}
\begin{enumerate}[leftmargin=2em,label=\arabic*.]
  \setcounter{enumi}{9}
  \item \textbf{Projection Logic and Interpretive Discipline}
    \begin{enumerate}[label*=\arabic*., leftmargin=2.5em]
      \item From structure to projection
      \item What is forced, what is read, what is hypothesized
      \item Dirac / QFT / Einstein-type readings
      \item Higgs as coupling instance or mean coupling point
    \end{enumerate}
\end{enumerate}

\paragraph{Part V --- Consolidation}
\begin{enumerate}[leftmargin=2em,label=\arabic*.]
  \setcounter{enumi}{10}
  \item \textbf{Consistency Layer}
    \begin{enumerate}[label*=\arabic*., leftmargin=2.5em]
      \item Internal structural consistency
      \item Probe/main compatibility
      \item Closure robustness
      \item Remaining soft spots
    \end{enumerate}
  \item \textbf{Outlook to v1.2.0}
    \begin{enumerate}[label*=\arabic*., leftmargin=2.5em]
      \item What comes next
      \item Planned lemma expansion
      \item Publication path
    \end{enumerate}
\end{enumerate}

\paragraph{Appendix scaffold}
\begin{enumerate}[leftmargin=2em,label=\Alph*.]
  \item Symbol table
  \item Structural tools catalog
  \item Technical lemmas
  \item Computational shortcuts
  \item Projection tables
  \item Probe sector schematics
\end{enumerate}

\subsection*{Minimum formal architecture reserved for v\docversion}
The next expansion pass should lock at least the following formal blocks:
\begin{itemize}[leftmargin=2em]
  \item Proposition: Ground Condition implies emergence of HC.
  \item Proposition: HC implies structural closure.
  \item Proposition: Closure confines the admissible dynamics to the septinion/seption regime.
  \item Proposition: The probe sector inherits the same closure discipline under orthogonal coupling.
  \item Proposition: The first-layer master equation is structurally forced by the closure architecture.
  \item Auxiliary lemmata: orthogonality, frame shift, confinement, probe compatibility, read-off consistency.
\end{itemize}

\subsection*{Practical build order}
For the actual writing sequence of the next cycle:
\begin{enumerate}[leftmargin=2em]
  \item Ground Condition / HC / Closure
  \item Probe sector
  \item Master equation (first layer)
  \item Tool normalization
  \item Interpretive bridge
  \item Consistency layer
  \item Appendix fill-in
\end{enumerate}


% ==================================================================
% v2.27.5 PROOF-CHAIN ROADMAP
% ==================================================================

\clearpage
